For succinctly described two-player games the minimax value has no invariant computable in deterministic polynomial time when play can be exponentially long; for play of polynomial length such an invariant exists exactly when $\Pclass=\PSPACE$; and an evaluator that sees payoffs only through queries needs exponentially many queries in the depth. We restate these classical limits in one model and then study what collapses exactly. The value is a signed spectral index of the unfolded game tree. For undirected vertex geography the index is a ratio of matching polynomials of the input graph, and for directed vertex geography it is a ratio of weighted matching polynomials of one fixed graph whenever every strong component is symmetric. We prove the converse for real weights: a digraph admits such a local weighted matching potential with weights in $\Qz$, in $\Rz$ or in real Puiseux series if and only if every strong component is symmetric. The proof applies Br\"and\'en's criterion for real stability to first order at $z=\infty$, and it also excludes Hermitian determinantal potentials and frozen environments. Over the algebraic closure of $\Qz$ the statement is false: a directed $m$-cycle together with $m+1$ isolated vertices carries a potential whose environment weights are the roots of an explicit polynomial, with full symmetric monodromy for $m\le6$. We also decide directed geography in polynomial time when the one-way arcs of every strong component share a tail and show that matching data decide no other two-arc geometry; prove that every two-periodic blow-up of a directed cycle carries a set potential; and show that exact resolvents on digraphs of small width need exponential degree and non-local memory, while circuits over component states have linear size.
